% =============================================================================
% main.tex
% The Joint Episode as an Ontological-Temporal Unit
% Technical-mathematical paper on the Joint Episode
% Initial skeleton — v0.1-initial (2026-06-18)
% =============================================================================

\documentclass[11pt,a4paper]{article}

\usepackage[margin=1.0in,top=0.9in,bottom=0.9in]{geometry}
\usepackage[T1]{fontenc}
\usepackage{lmodern}
\usepackage{setspace}
\onehalfspacing
\usepackage{microtype}

\usepackage{amsmath,amssymb,amsfonts}
\usepackage{booktabs}
\usepackage{enumitem}
\setlist[itemize]{leftmargin=*,itemsep=0.2em,topsep=0.15em}
\setlist[enumerate]{leftmargin=*,itemsep=0.2em,topsep=0.15em}

\usepackage{tcolorbox}
\tcbuselibrary{skins,breakable}
\newtcolorbox{conceptbox}[1][]{
    colback=blue!3!white,
    colframe=conceptcolor!80!black,
    boxrule=0.6pt,
    arc=2pt,
    left=6pt,
    right=6pt,
    top=3pt,
    bottom=3pt,
    fonttitle=\bfseries\small,
    title=#1,
    breakable
}
\newtcolorbox{importantbox}{
    colback=orange!4!white,
    colframe=orange!70!black,
    boxrule=0.6pt,
    arc=2pt,
    left=5pt,right=5pt,top=3pt,bottom=3pt
}

\usepackage{tikz}
\usetikzlibrary{shapes,arrows.meta,positioning,calc}

\usepackage{hyperref}
\hypersetup{
    colorlinks=true,
    linkcolor=black,
    urlcolor=blue,
    citecolor=black,
    pdfauthor={Dr. Johel Padilla Villaunueva},
    pdftitle={The Joint Episode as an Ontological-Temporal Unit: Mathematical Characterization, Structural Properties, and Transitions in Complex Systems},
    pdfsubject={Joint Episode; ontological-temporal unit; complex systems; recurrence quantification; systemic tau},
    pdfkeywords={joint episode; ontological layers; recurrence quantification analysis; systemic tau; antisynchrony; complex systems; structural transitions}
}

\usepackage{fancyhdr}
\pagestyle{fancy}
\fancyhf{}
\fancyhead[L]{\footnotesize The Joint Episode as an Ontological-Temporal Unit}
\fancyhead[R]{\footnotesize Working Draft}
\fancyfoot[C]{\footnotesize\thepage}
\renewcommand{\headrulewidth}{0.4pt}

\title{\textbf{The Joint Episode as an Ontological-Temporal Unit:}\\[0.25em]
\large Mathematical Characterization, Structural Properties, and Transitions in Complex Systems}

\author{
    \textbf{Dr. Johel Padilla Villaunueva}\\[0.2em]
    \small Ontología del Presente (Tau Sistémico) Research Program
}

\date{\today \\ \small Working Draft — Initial Structure}

\begin{document}

\maketitle

\begin{abstract}
% === ABSTRACT DRAFT TO BE REFINED ===
\noindent
The Joint Episode is proposed as a fundamental ontological-temporal unit in the analysis of complex systems. We provide a rigorous mathematical characterization grounded in recurrence quantification analysis (RQA) and the three-layer ontological framework of systemic persistence. A Joint Episode is formally defined as a maximal interval of sustained low antisynchrony (Layer 2 permeability) that contains at least one episode of local intensification (Layer 1: elevated critical mass $M$ derived from hyper-persistence of Systemic Tau $\tau_s$ and structural trapping time). 

We introduce and formalize the primary metrics: Joint Score (integrated intensity $\times$ duration), episode duration, mean/max intensity, and a calibrated Confidence Score. Structural properties are derived, including relations to anti-synchronization $A(k)$, active memory, and bidirectional probabilistic filtering between layers. Transitions are analyzed via layer dynamics in synthetic realizations of coupled logistic maps and the Lorenz attractor. The ontological implications are examined: the Joint Episode functions as a detectable \textit{kairos} — a coherent temporal window during which the probability of structural reorganization (Layer 3) is materially elevated by the conjunction of persistence intensification and scale coherence. 

This work establishes the Joint Episode as both a mathematically operationalizable construct and an ontologically primitive unit bridging local dynamics and global transitions in complex systems.

\vspace{0.15em}
\noindent\textbf{Keywords:} joint episode; ontological-temporal unit; systemic tau; recurrence quantification analysis; antisynchrony; complex systems; structural transitions; persistence; kairos.
\end{abstract}

\vspace{0.5em}

\section{Introduction}
% Content to be developed iteratively. See refined structure in project notes.

\section{Formal Definition of the Joint Episode}

\section{Metrics and Characterization}

\section{Structural Properties}

\section{Transitions and Layer Dynamics}

\section{Empirical Validation in Synthetic Systems}

\section{Ontological Implications}

\section{Discussion and Open Questions}

\section{Conclusions}

\bibliographystyle{plainnat}
\bibliography{references}

\end{document}
